%% kapa: exam=1 round=31 year=2020 period=2 form=A subject=appraisal_law subject_ko=감정평가관계법규 seq=9 number=9
%% source: Q-Net 감정평가사 기출문제 (cst00309) artlSeq=5207219 file=extracted/1008/5207219/A형/A형/2020년도 제31회 감정평가사 1차 시험_2교시_A형.pdf
%% conversion: kordoc 4.24.51 -> tools/gen1.py (auto); review=proofread-sonnet; flags=none
\begin{question}{9}
국토의 계획 및 이용에 관한 법령상 광역도시계획에 관한 설명으로 옳은 것은?
\begin{choices}
  \item 광역도시계획에는 경관계획에 관한 사항 중 광역계획권의 지정목적을 이루는 데 필요한 사항에 대한 정책 방향이 포함되어야 한다.
  \item 도지사가 광역계획권을 지정하려면 관계 중앙행정기관의 장의 의견을 들은 후 지방의회의 동의를 얻어야 한다.
  \item 광역도시계획을 공동으로 수립하는 시ㆍ도지사는 그 내용에 관하여 서로 협의가 되지 아니하는 경우 공동으로 국토교통부장관에게 조정을 신청하여야 한다.
  \item 광역계획권이 둘 이상의 시ㆍ도의 관할 구역에 걸쳐 있는 경우에는 국토교통부장관이 당해 광역도시계획의 수립권자가 된다.
  \item 도지사는 시장 또는 군수가 요청하는 경우에는 단독으로 광역도시계획을 수립할 수 있으며, 이 경우 국토교통부장관의 승인을 받아야 한다.
\end{choices}
\end{question}
\answer{1}
